\documentclass[10pt,a4paper]{article}
\usepackage[T1]{fontenc}
\usepackage[utf8]{inputenc}
\usepackage{lmodern}
\usepackage[margin=24mm]{geometry}
\usepackage{amsmath,amssymb,bm,booktabs,array}
\usepackage[hidelinks]{hyperref}
\usepackage{enumitem}
\newcommand{\CN}{\mathcal{CN}}
\newcommand{\E}{\mathbb{E}}
\newcommand{\I}{\mathbf{I}}
\newcommand{\vect}{\operatorname{vec}}
\newcommand{\tr}{\operatorname{tr}}
\newcommand{\diag}{\operatorname{diag}}
\newcommand{\logit}{\operatorname{logit}}
\newcommand{\sigmoid}{\operatorname{sigmoid}}
\newcommand{\clip}{\operatorname{clip}}
\newcommand{\relu}{\operatorname{Re}}
\newcommand{\imag}{\operatorname{Im}}
\newcommand{\norm}[1]{\left\lVert#1\right\rVert}
\renewcommand{\thesection}{\Roman{section}}
\renewcommand{\thesubsection}{\Alph{subsection}}
\setlength{\parskip}{3pt}
\setlength{\parindent}{1.2em}
\allowdisplaybreaks
\title{Joint Activity Detection and Channel Estimation\protect\\
for AFDM LEO Satellite Uplinks\protect\\[5pt]
\large Sections II and III}
\author{}
\date{Draft: October 8, 2026}
\begin{document}
\maketitle
\vspace{-7mm}
\setcounter{section}{1}

\section{System Model and Problem Formulation}
\label{sec:system}

\subsection{Grant-Free Access and Array Channel}

Consider the uplink of a low Earth orbit (LEO) satellite supporting grant-free
access from $K$ single-antenna terminals. The satellite employs a uniform
planar array (UPA) with $M=M_yM_z$ antennas. Each access block contains $Q$
affine frequency division multiplexing (AFDM) pilot frames, each with $N$
useful samples. The channel geometry and path coefficients remain fixed
within the access block. Terminal $k$ has activity indicator
$\alpha_k\in\{0,1\}$, and the receiver jointly estimates the terminal
activities and their effective channels.

Let $p$ be the terminal access probability before conditioning on a nonempty
access block. The activity law used to generate such blocks is
\begin{equation}
 p_{\rm gen}(\bm\alpha)=
 \frac{\mathbf{1}\{\sum_{k=1}^{K}\alpha_k\geq1\}}{Z}
 \prod_{k=1}^{K}p^{\alpha_k}(1-p)^{1-\alpha_k},
 \qquad Z=1-(1-p)^K.
 \label{eq:activity-generation}
\end{equation}
Its single-terminal marginal is $\rho=p/Z$. This publicly specified
marginal, together with the observations, determines the activity prior
used by the receiver.

We fix the complex-envelope convention by
$s_{\rm RF}(t)=\relu\{x(t)e^{-j2\pi f_ct}\}$ and downconversion by
$e^{+j2\pi f_ct}$. A path with time-varying propagation delay $\tau(t)$
then contributes $g\,x(t-\tau(t))e^{+j2\pi f_c\tau(t)}$. Static carrier
phase is absorbed into the path gain. Under the narrowband block model,
envelope delays and array directions are held fixed while the first-order
carrier phase variation is retained.

Denote the satellite position and velocity by $\mathbf{s}$ and
$\mathbf{v}_{\rm sat}$, respectively, and the position of terminal $k$ by
$\mathbf{g}_k$. With propagation speed $c_0$ and carrier frequency $f_c$,
the line-of-sight direction and common terminal Doppler shift are
\begin{equation}
 \mathbf{u}_k=\frac{\mathbf{g}_k-\mathbf{s}}
 {\norm{\mathbf{g}_k-\mathbf{s}}},
 \qquad \nu_k=-\frac{f_c}{c_0}\mathbf{v}_{\rm sat}^{\mathsf T}\mathbf{u}_k.
 \label{eq:geometry}
\end{equation}
Indeed, the range $L_k=\norm{\mathbf g_k-\mathbf s}$ obeys
$\dot L_k=-\mathbf v_{\rm sat}^{\mathsf T}\mathbf u_k$, so the chosen
envelope has Doppler $\nu_k=f_c\dot L_k/c_0$ and time factor
$e^{+j2\pi\nu_kt}$. The common reference delay below is fixed at the
block reference time; it does not additionally remove a common Doppler.

Let $\tau_c$ be the propagation time from the satellite to the closest
ground point in the beam footprint. After compensation of this common
reference delay, path $\ell$ has residual delay
\begin{equation}
 \tau_{k\ell}=\frac{\norm{\mathbf{g}_k-\mathbf{s}}}{c_0}-\tau_c
             +\delta_\ell\tau_{{\rm spread},k},
 \label{eq:path-delay}
\end{equation}
where $\delta_\ell$ and $\tau_{{\rm spread},k}$ specify the relative path
delay and terminal delay spread. All paths of a terminal share its spatial
direction and Doppler shift, following the common spatial response model
for satellite links in \cite[Sec.~II-C]{li2023} and the common Doppler
model in \cite[Sec.~II]{huang2026}.

Let $\mu_{y,k}$ and $\mu_{z,k}$ be the projections of $\mathbf{u}_k$ onto
the fixed unit axes $\mathbf e_y,\mathbf e_z$. The array reference and
indexing convention are
$\mathbf s_m=\mathbf s-n_y d_y\mathbf e_y-n_zd_z\mathbf e_z$.
The far-field delay difference is
$(n_y d_y\mu_{y,k}+n_zd_z\mu_{z,k})/c_0$; with the envelope convention
above, the narrowband array response at antenna index $m=n_z+M_zn_y$ is
\begin{equation}
 [\mathbf{a}_k]_m=
 \exp\!\left[j2\pi\left(
 \frac{d_y}{\lambda}n_y\mu_{y,k}
 +\frac{d_z}{\lambda}n_z\mu_{z,k}\right)\right],
 \qquad \norm{\mathbf{a}_k}^2=M,
 \label{eq:upa}
\end{equation}
where $\lambda=c_0/f_c$. The unknown spatial response is included in the
channel to be estimated.

\subsection{AFDM Channel with Continuous Delays}

Let $B=N\Delta f$ be the sampling rate, $T_s=1/B$, and $N_{\rm CPP}$ the
chirp-periodic prefix (CPP) length. Following the AFDM modulation in
\cite[Secs.~II-B and III-A]{bemani2022}, the discrete affine Fourier
transform (DAFT) matrix $\mathbf{F}_{\rm A}$ is defined by
\begin{equation}
 [\mathbf{F}_{\rm A}]_{rt}=\frac{1}{\sqrt N}
 \exp[-j2\pi(c_2r^2+rt/N+c_1t^2)],
 \quad r,t\in\{0,\ldots,N-1\}.
 \label{eq:daft}
\end{equation}
An AFDM pilot vector $\mathbf{p}_{qk}$ is mapped by the inverse discrete
affine Fourier transform (IDAFT) to
$\mathbf{s}_{qk}=\mathbf{F}_{\rm A}^{\mathsf H}\mathbf{p}_{qk}$. The
prefix extends the integer-indexed samples according to
\begin{equation}
 s_{qk}[t]=\exp\{j2\pi c_1(t^2-\langle t\rangle_N^2)\}
              s_{qk}[\langle t\rangle_N],
 \quad -N_{\rm CPP}\leq t<0,
 \label{eq:cpp}
\end{equation}
where $\langle t\rangle_N=t-N\lfloor t/N\rfloor$. The prefix covers the
residual channel delay, $N_{\rm CPP}T_s\geq\max_{k,\ell}\tau_{k\ell}$.
This expression retains the chirp phase over every prefix period. The frame duration is
$T_{\rm frame}=(N+N_{\rm CPP})T_s$.

To define fractional delays, we specify the continuous waveform between
the discrete samples. For real normalized time $\xi$, set
$w_\xi=\langle\xi\rangle_N$, $\zeta_n(w)=\lfloor2c_1w+n/N\rfloor$, and
\begin{equation*}
 \varphi_n(\xi)=\frac{1}{\sqrt N}
 \exp\!\left[j2\pi\left(c_2n^2+c_1\xi^2+\frac{nw_\xi}{N}
 -\zeta_n(w_\xi)w_\xi\right)\right].
 \tag{6a}\label{eq:continuous-waveform}
\end{equation*}
This is the adopted wrapped-chirp waveform model, rather than a
consequence of discrete samples alone. It agrees with the IDAFT at
integer useful-sample times and obeys the same cross-period chirp
extension as \eqref{eq:cpp}. For a delay of $d=\tau/T_s$ samples,
define $\eta=t-d$, $w=\langle\eta\rangle_N$, and
$\zeta_{tn}=\lfloor2c_1w+n/N\rfloor$. Substitution into
$\varphi_n(t-d)$ gives
\begin{equation}
 [\mathbf{B}_d]_{tn}=\frac{1}{\sqrt N}
 \exp\!\left[j2\pi\left(c_2n^2+c_1\eta^2
                  +\frac{nw}{N}-\zeta_{tn}w\right)\right].
 \label{eq:continuous-basis}
\end{equation}
We use the Doppler modulation convention $\exp(+j2\pi\nu tT_s)$.
For useful-window start index $g$, the corresponding path matrix is
\begin{equation}
 \mathbf{G}(d,\nu;g)=\mathbf{F}_{\rm A}
 \diag\{e^{j2\pi\nu T_s(g+t)}\}_{t=0}^{N-1}\mathbf{B}_d.
 \label{eq:path-operator}
\end{equation}
Expanding its entries, with $2Nc_1\in\mathbb N_+$ and integer receive
index $t$, gives
\begin{equation*}
 \begin{split}
 G_{rn}&=\frac{e^{j2\pi[c_1d^2+c_2(n^2-r^2)-nd/N+\nu T_sg]}}{N}\\
 &\quad\times\sum_{t=0}^{N-1}
 e^{j2\pi[(n-r)/N-2c_1d+\nu T_s]t}
 e^{j2\pi\zeta_{tn}d}.
 \end{split}
 \tag{8a}\label{eq:path-entry-expansion}
\end{equation*}
The interval-wise sum is the fractional-delay form in
\cite[Sec.~II, eq.~(5)]{huang2026}, with Doppler signs explicitly
converted: $k_{\nu,\mathrm{Huang}}=-\nu NT_s$. The useful-window phase
is retained in \eqref{eq:path-operator} and can be included in the
reference-frame path gain. At integer $d$, the final exponential in
\eqref{eq:path-entry-expansion} equals one.

The chirp parameters satisfy $2Nc_1\in\mathbb{N}_{+}$. At an integer
delay $d$, \eqref{eq:path-operator} reduces to
\begin{align}
 \mathbf{G}(d,\nu;g)&=\mathbf{F}_{\rm A}\mathbf{D}_{\nu,g}
             \bm\Omega_d\bm\Pi_{d\bmod N}\mathbf{F}_{\rm A}^{\mathsf H},
 \label{eq:integer-operator}\\
 [\bm\Omega_d]_{tt}&=
 e^{j2\pi c_1[(t-d)^2-\langle t-d\rangle_N^2]},
 \label{eq:integer-prefix-phase}
\end{align}
where $\bm\Pi_l\mathbf{s}$ has entries $s[\langle t-l\rangle_N]$.
Thus the full residual delay enters the prefix phase, while its remainder
determines the cyclic shift.

At the reference frame, whose useful-window start is $N_{\rm CPP}$,
terminal $k$ has the DAFT-domain channel
\begin{equation}
 \mathbf{G}_k=\sum_{\ell}g_{k\ell}
 \mathbf{G}(\tau_{k\ell}/T_s,\nu_k;N_{\rm CPP}),
 \qquad H_{kmrn}=[\mathbf{a}_k]_m[\mathbf{G}_k]_{rn}.
 \label{eq:channel}
\end{equation}
Here $g_{k\ell}$ includes the complex path gain after large-scale power
compensation. The tensor $H_k\in\mathbb{C}^{M\times N\times N}$ contains
all antenna, DAFT output, and DAFT input coordinates. Its spatial unfolding
$\mathbf{H}_k\in\mathbb{C}^{M\times N^2}$ has rank at most one under the
common spatial response assumption. The receiver learns the distribution
of this full tensor from channel samples.

\subsection{Pilots and the Received Signal}

We assign pilots using randomly selected rows and permuted columns of a
$KN$-point discrete Fourier transform (DFT). With distinct row indices
$r_q$, column permutation $c_{kn}$, and independent uniform phases
$\phi_{kn}$, the pilots are
\begin{equation}
 P_{qkn}=\sqrt{\frac{E_p}{N}}
 e^{-j2\pi r_qc_{kn}/(KN)}e^{j\phi_{kn}},
 \qquad \sum_{n=0}^{N-1}|P_{qkn}|^2=E_p.
 \label{eq:pilot}
\end{equation}
Because the terminal Doppler is common to all of its paths, moving the
useful window by one frame multiplies its channel by $e^{j\beta_k}$,
where $\beta_k=\operatorname{wrap}(2\pi\nu_kT_{\rm frame})$. We assume
that these effective cross-frame phases are available to the receiver and
absorb them into the pilot coefficients:
\begin{equation}
 A_{qkn}=P_{qkn}e^{jq\beta_k},\qquad q=0,\ldots,Q-1.
 \label{eq:effective-pilot}
\end{equation}
The general matrix $\mathbf{A}$ defined below is used in inference,
including the terminal-dependent phases.

Define the activity-weighted channel $X_k=\alpha_kH_k$. After prefix
removal and DAFT demodulation, the observation is
\begin{equation}
 Y_{qmr}=\sum_{k=1}^{K}\sum_{n=0}^{N-1}A_{qkn}X_{kmrn}+W_{qmr},
 \qquad W_{qmr}\sim\CN(0,\sigma^2).
 \label{eq:observation-tensor}
\end{equation}
Noise entries are conditionally independent. To expose the shared linear
operator, set
\begin{equation}
 [\mathbf{Y}]_{q,mN+r}=Y_{qmr},\quad
 [\mathbf{A}]_{q,(k-1)N+n}=A_{qkn},\quad
 [\mathbf{X}]_{(k-1)N+n,mN+r}=X_{kmrn}.
 \label{eq:mmv-indexing}
\end{equation}
Then
\begin{equation}
 \mathbf{Y}=\mathbf{A}\mathbf{X}+\mathbf{W},\quad
 \mathbf{Y}\in\mathbb{C}^{Q\times MN},\quad
 \mathbf{A}\in\mathbb{C}^{Q\times KN},\quad
 \mathbf{X}\in\mathbb{C}^{KN\times MN}.
 \label{eq:mmv}
\end{equation}
This is a multiple measurement vector (MMV) model with $MN$ measurement
vectors and a shared terminal activity pattern. Equivalently,
$\mathbf{y}=\mathcal{A}\mathbf{x}+\mathbf{w}$, with
$\mathcal{A}=\I_{MN}\otimes\mathbf{A}$ and
$\mathbf{x}=\vect(\mathbf{X})$. Grouping the coordinates of each terminal
uses the fixed permutation specified by \eqref{eq:mmv-indexing}.

The receiver is supplied with $(\mathbf{Y},\mathbf{A},\sigma^2,\rho)$.
In the simulation model, the received signal-to-noise ratio (SNR) is set
through the clean superposition $\mathbf{Y}_0=\mathbf{A}\mathbf{X}$:
\begin{equation}
 \sigma^2=\frac{\norm{\mathbf{Y}_0}_F^2}{QMN}
                     10^{-\mathrm{SNR}_{\rm dB}/10}.
 \label{eq:snr}
\end{equation}
The supplied noise variance is the parameter of the complex Gaussian
noise law. The inference variables are expressed in fixed training units:
$\mathbf{Y}\leftarrow\mathbf{Y}/s$,
$H_k\leftarrow H_k/s$, $X_k\leftarrow X_k/s$,
$\mathbf{W}\leftarrow\mathbf{W}/s$, and $\sigma^2\leftarrow\sigma^2/s^2$, with a
common scale $s$. We retain the same symbols for the normalized variables.

\subsection{Bayesian Formulation}

Let $D=MN^2$ be the number of complex coordinates per terminal, and let
$q_{\rm tr}$ denote the normalized candidate-channel distribution used
for training. It is obtained by drawing terminals from the full geometry
model before applying their activity masks. The receiver uses the
factorized prior
\begin{equation}
 q(\{dH_k\},\bm\alpha)=
 \prod_{k=1}^{K}q_{\rm tr}(dH_k)
       \rho^{\alpha_k}(1-\rho)^{1-\alpha_k},
 \label{eq:working-prior}
\end{equation}
and forms the working posterior
\begin{equation}
 \begin{aligned}
 q(d\mathbf{X},\{dH_k\},\bm\alpha\mid\mathbf{Y},\mathbf{A},\sigma^2)
 &\propto e^{-\norm{\mathbf{Y}-\mathbf{A}\mathbf{X}}_F^2/\sigma^2}\\
 &\quad\times\prod_{k=1}^{K}\delta_{\alpha_kH_k}(dX_k)
       q_{\rm tr}(dH_k)\rho^{\alpha_k}(1-\rho)^{1-\alpha_k}.
 \end{aligned}
 \label{eq:working-posterior}
\end{equation}
Here $\delta_{\alpha_kH_k}$ enforces $X_k=\alpha_kH_k$. The strict
physical prior conditioned on the supplied quantities is
$\Pi_{\rm phys}(\{dH_k\},\bm\alpha\mid\mathbf A,\sigma^2)$:
the operator shares geometric variables with the channels, and
\eqref{eq:snr} makes the noise variance depend on the clean received
power. It also contains the nonempty access constraint.
Equation~\eqref{eq:working-prior} replaces this conditional joint law by
the candidate-channel marginal and independent activity marginals.
This is a working-model approximation; the common Gaussian likelihood
does not establish equality of the two posteriors.

The desired outputs are the effective channel estimate and activity
probability for each terminal. Under the working posterior, their Bayesian
targets are
\begin{equation}
 \widehat X_k=\E_q[X_k\mid\mathbf{Y},\mathbf{A},\sigma^2],\qquad
 p_k=q(\alpha_k=1\mid\mathbf{Y},\mathbf{A},\sigma^2),\qquad
 \widehat\alpha_k=\mathbf{1}\{p_k\geq\tfrac12\}.
 \label{eq:targets}
\end{equation}
The following receiver approximates these quantities by exchanging
Gaussian messages between a linear estimator and a learned channel
denoiser, with a classifier providing the activity density ratio.

\section{Receiver Algorithm}
\label{sec:algorithm}

The receiver follows the linear-estimation and nonlinear-denoising
decomposition of score-based turbo message passing
\cite[Secs.~III--V]{cai2026}. The linear module incorporates
$\mathbf{A}$ and $\sigma^2$. The nonlinear module processes the complete
terminal channel tensor and returns its channel and activity moments.
A first network provides the channel mean, a second network learns its
conditional error, and an activity classifier estimates the ratio between
active-channel and zero-signal query densities. Gaussian message division
and moment projection connect the modules
\cite[Secs.~II--III]{minka2001}.

\subsection{Linear Estimation and Extrinsic Queries}

Suppose the incoming message has mean
$\mathbf{X}^{\rm pri}\in\mathbb{C}^{KN\times MN}$ and common covariance
$\gamma\I_{KN}$ for each measurement vector, with $\gamma>0$,
$\sigma^2>0$, and $\mathbf A\ne\mathbf0$. Gaussian conditioning gives
\begin{align}
 \mathbf{X}_{\rm L}
 &=\mathbf{X}^{\rm pri}+\gamma\mathbf{A}^{\mathsf H}
   (\sigma^2\I_Q+\gamma\mathbf{A}\mathbf{A}^{\mathsf H})^{-1}
   (\mathbf{Y}-\mathbf{A}\mathbf{X}^{\rm pri}),
 \label{eq:linear-mean}\\
 \bm\Sigma_{\rm L}
 &=(\gamma^{-1}\I_{KN}+\sigma^{-2}\mathbf{A}^{\mathsf H}\mathbf{A})^{-1},
 \qquad t_{\rm L}=\frac{\tr\bm\Sigma_{\rm L}}{KN}.
 \label{eq:linear-variance}
\end{align}
The covariance eigenvalues imply $0<t_{\rm L}<\gamma$. Projecting the
covariance onto $t_{\rm L}\I_{KN}$ and dividing out the incoming
Gaussian message therefore yields positive, finite extrinsic variance:
\begin{equation}
 v_{\rm L}=(t_{\rm L}^{-1}-\gamma^{-1})^{-1},\qquad
 \mathbf{R}_{\rm L}=v_{\rm L}
     \left(\frac{\mathbf{X}_{\rm L}}{t_{\rm L}}
                         -\frac{\mathbf{X}^{\rm pri}}{\gamma}\right).
 \label{eq:linear-extrinsic}
\end{equation}
The rows of $\mathbf{R}_{\rm L}$ are reorganized into terminal tensors
$R_k\in\mathbb{C}^{M\times N\times N}$ using
\eqref{eq:mmv-indexing}. Damping of these parameters is specified in
Section~\ref{sec:schedule}.

Efficient evaluation uses a thin singular value decomposition (SVD)
$\mathbf{A}=\mathbf{U}\diag(s_i)\mathbf{V}^{\mathsf H}$. Let
\begin{equation}
 f_{\rm L}=\frac{1}{KN}\sum_i
       \frac{\gamma s_i^2}{\sigma^2+\gamma s_i^2}.
 \label{eq:linear-fraction}
\end{equation}
Then $t_{\rm L}=\gamma(1-f_{\rm L})$,
$v_{\rm L}=t_{\rm L}/f_{\rm L}$, and
$\mathbf{R}_{\rm L}=\mathbf{X}^{\rm pri}
+(\mathbf{X}_{\rm L}-\mathbf{X}^{\rm pri})/f_{\rm L}$.
The decomposition is computed once per observation operator.

\subsection{Learning the Channel Mean and Conditional Error}

The learned denoiser is trained on the pseudo additive white Gaussian
noise (AWGN) model
\begin{equation}
 R=H+W,\qquad H\sim q_{\rm tr},\qquad
 \vect(W)\sim\CN(\mathbf{0},v\I_D).
 \label{eq:pseudo-awgn}
\end{equation}
Equation~\eqref{eq:pseudo-awgn} is exact for the generated training
queries and a surrogate for the finite-dimensional AFDM receiver queries.
To see the distinction, define
$\mathbf L=\gamma\mathbf A^{\mathsf H}
(\sigma^2\I+\gamma\mathbf A\mathbf A^{\mathsf H})^{-1}/f_{\rm L}$.
The undamped linear query obeys the identity
\begin{equation*}
 \mathbf R_{\rm L}-\mathbf X
 =(\mathbf L\mathbf A-\I)(\mathbf X-\mathbf X^{\rm pri})
 +\mathbf L\mathbf W.
 \tag{25a}\label{eq:linear-query-error}
\end{equation*}
This identity does not establish that the actual error is independent,
white, or Gaussian. The scalar AWGN interpretation and damping are used
with this approximation explicit, as discussed for finite systems in
\cite[Sec.~V-B1]{cai2026}.

Write
$\mathbf{u}=(\relu\vect(R),\imag\vect(R))\in\mathbb{R}^{2D}$,
define $\mathbf{h}_{\mathbb R}$ similarly, and let $\nu=v/2$.
We assume the moments needed by the losses are finite; finite fourth
moments and the corresponding network integrability suffice for the
conditional squared-error calculations below. Positive-variance Gaussian
convolution has a smooth positive density even for a channel prior
supported on a lower-dimensional set. If $q_v^{\mathbb R}$ is this
active query density in real coordinates, Tweedie's identities
\cite{efron2011,cai2026} give
\begin{align}
 \mathbf{m}_{\mathbb R}(\mathbf{u},v)
 &=\mathbf{u}+\nu\nabla_{\mathbf{u}}\log q_v^{\mathbb R}(\mathbf{u}),
 \label{eq:tweedie-mean}\\
 \mathbf{C}_{\mathbb R}(\mathbf{u},v)
 &=\nu\I_{2D}+\nu^2\nabla_{\mathbf{u}}^2\log q_v^{\mathbb R}(\mathbf{u}).
 \label{eq:tweedie-covariance}
\end{align}
We train $S_\theta(\mathbf{u},v)\in\mathbb{R}^{2D}$ using denoising
score matching \cite{vincent2011,cai2026}. For
$\bm\epsilon=(\mathbf{u}-\mathbf{h}_{\mathbb R})/\sqrt\nu$, its objective is
\begin{equation}
 \mathcal{L}_1(\theta)=\E\left[
 \frac{1}{2D}\norm{\sqrt\nu S_\theta(\mathbf{u},v)+\bm\epsilon}^2\right].
 \label{eq:first-loss}
\end{equation}
The resulting full channel mean is
\begin{equation}
 F_\theta(R,v)=R+\frac v2
       \left(S_{\theta,\Re}(R,v)+jS_{\theta,\Im}(R,v)\right).
 \label{eq:learned-mean}
\end{equation}
All $M N^2$ complex coordinates are retained in this mapping.

The second network $T_\psi(\mathbf{u},v)\in\mathbb{R}^{2D}$ is trained
after freezing $S_\theta$. Define the scalar readout
\begin{equation}
 a_\psi(R,v)=1+\nu\frac{1}{2D}\sum_{i=1}^{2D}T_{\psi,i}(\mathbf{u},v),
 \qquad \widehat\tau_\psi(R,v)=v a_\psi(R,v).
 \label{eq:second-readout}
\end{equation}
With $\mathbf{b}=\sqrt\nu S_\theta+\bm\epsilon$, the global-average
second-order training objective is
\begin{align}
 \mathcal{L}_2(\psi)
 &=\E\left[\left\{
 \frac{1}{2D}\sum_{i=1}^{2D}(\nu T_{\psi,i}-b_i^2+1)
                            \right\}^2\right]\notag\\
 &=\E\left[\left\{
 a_\psi(R,v)-\frac{\norm{F_\theta(R,v)-H}_F^2}{Dv}
                            \right\}^2\right].
 \label{eq:second-loss}
\end{align}
The second equality follows from
$F_{\theta,\mathbb R}-\mathbf{h}_{\mathbb R}=\sqrt\nu\mathbf{b}$ and
$2D\nu=Dv$. It identifies the statistic learned by the network: at the
function-space optimum, the scalar variance readout is
\begin{equation}
 \tau_\theta(R,v)=\frac{1}{D}\E[
 \norm{F_\theta(R,v)-H}_F^2\mid R,v,\alpha=1].
 \label{eq:conditional-risk}
\end{equation}
Writing $m(R,v)=\E[H\mid R,v,\alpha=1]$ and
$\tau_{\rm C}=D^{-1}\E[\norm{H-m}_F^2\mid R,v,\alpha=1]$, conditional
orthogonality gives
\begin{equation}
 \tau_\theta(R,v)=\tau_{\rm C}(R,v)
             +\frac{\norm{F_\theta(R,v)-m(R,v)}_F^2}{D}.
 \label{eq:risk-decomposition}
\end{equation}
Thus the learned error includes the approximation error of the deployed
mean. When $F_\theta=m$, it reduces to the posterior covariance trace
per complex coordinate. For $U=\norm{F_\theta-H}_F^2/(Dv)$, the loss
decomposes as
$\E[(a_\psi-\E[U\mid R,v,\alpha=1])^2]
+\E[\operatorname{Var}(U\mid R,v,\alpha=1)]$.
Only the complete trace is supervised: any perturbation of $T_\psi$
whose coordinates sum to zero leaves \eqref{eq:second-loss} unchanged.
Accordingly, individual Hessian diagonal entries or per-antenna
subtraces are not separately identified by this objective.

For a receiver query $(\{R_k\},v)$, we compute
$F_k=F_\theta(R_k,v)$ and $\widehat\tau_k=\widehat\tau_\psi(R_k,v)$,
and use the common scalar
\begin{equation}
 \tau_{\rm raw}=\frac1K\sum_{k=1}^{K}\widehat\tau_k,\qquad
 \tau=\max\{\tau_{\rm raw},\tau_{\min}\}.
 \label{eq:pooled-variance}
\end{equation}
This averaging fits a common isotropic Gaussian around the fixed means
$F_k$, giving equal weight to each terminal's active-conditional
distribution. The conditional expected negative log density, up to a
constant, is $KD\log\tau+(D/\tau)\sum_k\tau_\theta(R_k,v)$.
Its derivative is $KD(\tau-K^{-1}\sum_k\tau_\theta)/\tau^2$; hence the
positive-risk minimizer is the average conditional error. With the
constraint $\tau\geq\tau_{\min}>0$, the minimizer is its maximum with
the floor, also covering zero risk. Equation~\eqref{eq:pooled-variance}
substitutes finite network readouts for these conditional risks.

\subsection{Activity Detection by Density Ratio Estimation}

The activity classifier compares queries generated by an active channel
with queries containing only noise. Their densities under
\eqref{eq:pseudo-awgn} are
\begin{equation}
 \begin{aligned}
 q_{1,v}(R)&=\int\CN(\vect(R);\vect(H),v\I_D)q_{\rm tr}(dH),\\
 q_{0,v}(R)&=\CN(\vect(R);\mathbf{0},v\I_D).
 \end{aligned}
 \label{eq:query-densities}
\end{equation}
Both classes use the same training distribution of $v$. With activity
label $a\in\{0,1\}$ and training class prior $\pi_{\rm tr}=1/2$, the
classifier logit $z_\phi(R,v,F_\theta(R,v))$ minimizes binary cross
entropy (BCE):
\begin{equation}
 \mathcal{L}_{\rm act}(\phi)=\E[
 \log(1+e^{z_\phi})-a z_\phi].
 \label{eq:classifier-loss}
\end{equation}
The local posterior underlying these targets is proportional to
$\rho^a(1-\rho)^{1-a}\CN(R;aH,v\I_D)q_{\rm tr}(dH)$, with normalizer
$(1-\rho)q_{0,v}(R)+\rho q_{1,v}(R)$. This connects channel denoising
and activity inference within the same surrogate query model.

The channel mean network is frozen when its output is used as a feature.
Writing $\eta_{\rm tr}=\Pr_{\rm tr}(a=1\mid R,v)$, the conditional BCE
derivative is $\sigmoid(z)-\eta_{\rm tr}$ and its second derivative is
positive. The common variance-sampling density cancels in the posterior
odds. Thus, at the function-space optimum,
\begin{equation}
 z^*(R,v)=\logit\pi_{\rm tr}
                       +\log\frac{q_{1,v}(R)}{q_{0,v}(R)}.
 \label{eq:ideal-logit}
\end{equation}
This is the standard relationship between class-probability estimation
and density ratio estimation \cite[Secs.~III--IV]{menon2016}.
Keeping the complete $R$ in the input preserves the query information,
while $F_\theta(R,v)$ provides a deterministic feature for the finite
classifier. The BCE excess risk is the expected Kullback--Leibler (KL) divergence between
the true and estimated training conditional Bernoulli laws. This
population statement does not establish calibration of the finite
classifier on actual iterative queries.

The receiver corrects for the training class prior and incorporates the
terminal access marginal once:
\begin{equation}
 \ell_k=z_\phi(R_k,v,F_k)-\logit\pi_{\rm tr},\qquad
 p_k=\sigmoid(\logit\rho+\ell_k),
 \label{eq:activity-probability}
\end{equation}
where $\logit a=\log[a/(1-a)]$ and
$\sigmoid b=(1+e^{-b})^{-1}$. Under unchanged class-conditional query
laws, changes in $\rho$ are handled by this explicit prior-odds correction.

For implementation, the classifier receives the real and imaginary maps
of $R$ and $F_\theta$, scaled by $\sqrt{2/(1+v)}$, together with
\begin{equation}
 \begin{aligned}
 \Big(&\log v/5,\ \log(1+e_R),\ \log(1+e_F),\\
      &\operatorname{sgn}(c_{RF})\log(1+|c_{RF}|),\
       \log(1+e_{\rm res})\Big),
 \end{aligned}
 \label{eq:classifier-features}
\end{equation}
where $e_R=\norm{R}_F^2/(Dv)$,
$e_F=\norm{F_\theta}_F^2/(Dv)$,
$c_{RF}=\relu\langle R,F_\theta\rangle_F/(Dv)$, and
$e_{\rm res}=\norm{R-F_\theta}_F^2/(Dv)$. The feature maps have $4M$
channels. Three convolutional blocks with widths $64$, $128$, and $128$
are conditioned on $\log v/5$; a scalar readout produces $z_\phi$.

\subsection{Conditional Gaussian Messages and Joint Moments}

For each terminal we retain a Gaussian site with mean $E_k$ and scalar
variance $c_k$. When the query variance is within the trained range and
$0<\tau<v$, Gaussian division gives the updated site
\begin{equation}
 c_k=(\tau^{-1}-v^{-1})^{-1},\qquad
 E_k=c_k\left(\frac{F_k}{\tau}-\frac{R_k}{v}\right).
 \label{eq:site-update}
\end{equation}
The inequality makes the site precision positive; it is an acceptance
condition, not a general pointwise bound on the exact posterior variance.
For finite valid queries, an unsupported variance or $\tau\geq v$
retains the previous site. Invalid covariance or nonfinite network output
is reported as a failed run. Combining the available site with the
current query produces active-channel moments
\begin{equation}
 t_k=\frac{c_kv}{c_k+v},\qquad
 \mu_k=\frac{c_kR_k+vE_k}{c_k+v}.
 \label{eq:conditional-moments}
\end{equation}
An accepted update satisfies $\mu_k=F_k$ and $t_k=\tau$.
Within the trained range, activity detection in
\eqref{eq:activity-probability} uses the current $R_k,v,F_k$ whether
the site is updated or retained.

Outside that range, activity is computed from the normalized Gaussian
site densities:
\begin{align}
 \ell_k^{\rm G}
 &=\frac{\norm{R_k}_F^2}{v}
   -\frac{\norm{R_k-E_k}_F^2}{c_k+v}
   -D\log\!\left(1+\frac{c_k}{v}\right),
 \label{eq:gaussian-evidence}\\
 p_k&=\sigmoid(\logit\rho+\ell_k^{\rm G}).\notag
\end{align}
This follows by integrating the product
$\CN(R;H,v\I_D)\CN(H;E_k,c_k\I_D)$, which yields
$\CN(R;E_k,(v+c_k)\I_D)$, and dividing by the inactive density.
It is also the analytical activity rule used in the Gaussian-evidence
ablation \cite[Sec.~III-B2]{cai2026}. Since the site parameters depend
on the queries, this exact Gaussian integral is not an exact evidence
identity for the original physical prior.

The activity probability and conditional Gaussian moments define the
terminal approximation
\begin{equation}
 \widetilde q_k(dX_k)=(1-p_k)\delta_0(dX_k)
             +p_k\CN(\vect(X_k);\vect(\mu_k),t_k\I_D)\,dX_k.
 \label{eq:mixture}
\end{equation}
Consequently, the joint mean is $\widehat X_k=p_k\mu_k$. Let
$\mathbf{U}_k\in\mathbb{C}^{M\times N^2}$ be the spatial unfolding of
$\mu_k$. The spatial covariance averaged over the $N^2$ channel columns
is
\begin{equation}
 \mathbf{C}_{X,k}=p_kt_k\I_M
              +\frac{p_k(1-p_k)}{N^2}\mathbf{U}_k\mathbf{U}_k^{\mathsf H}.
 \label{eq:joint-spatial-covariance}
\end{equation}
These moments follow from total expectation and total covariance under
\eqref{eq:mixture}. The first term describes active-channel uncertainty,
and the second accounts for uncertainty in terminal activity. The finite
classifier and channel networks define the selected approximation; they
are not asserted to be the exact posterior moments of one fixed
generative prior. The common variance sent to the linear module is the
average trace
\begin{equation}
 \overline t_X=\frac{1}{KM}\sum_{k=1}^{K}\tr\mathbf{C}_{X,k}
 =\frac1K\sum_{k=1}^{K}\left[
 p_kt_k+p_k(1-p_k)\frac{\norm{\mu_k}_F^2}{D}\right].
 \label{eq:joint-average-variance}
\end{equation}
Define $b=\clip(\overline t_X/v,\varepsilon,1-\varepsilon)$. The outgoing
Gaussian message is
\begin{equation}
 \gamma_{\rm out}=\frac{vb}{1-b},\qquad
 X_{k,\rm out}=\frac{\widehat X_k-bR_k}{1-b}.
 \label{eq:nonlinear-extrinsic}
\end{equation}
The clipping keeps the outgoing precision positive and finite.
Equation~\eqref{eq:nonlinear-extrinsic} divides the adjusted Gaussian
with mean $\widehat X$ and variance $bv$ by the incoming Gaussian;
when clipping is active, $bv$ differs from the raw mixture variance.

\subsection{Iteration Schedule and Implementation}
\label{sec:schedule}

Initialize the linear prior mean at zero and set
$\gamma^{(0)}=\rho\widehat P_H$, where
$\widehat P_H$ is the training average of $\norm{H}_F^2/D$.
Initialize $E_k=0$, $c_k=1$, and an activity history $w_k=\rho$.
At iteration $i$, the linear module first computes
$(R_{{\rm L},k}^{(i)},v_{\rm L}^{(i)})$. From the second iteration onward,
the nonlinear query is damped as
\begin{equation}
 R_k^{(i)}=(1-\eta)R_k^{(i-1)}+\eta R_{{\rm L},k}^{(i)},\qquad
 v^{(i)}=(1-\eta)v^{(i-1)}+\eta v_{\rm L}^{(i)}.
 \label{eq:input-damping}
\end{equation}
After nonlinear estimation, the outgoing parameters are damped similarly:
\begin{equation}
 X_k^{{\rm pri},(i)}=(1-\eta)X_k^{{\rm pri},(i-1)}+\eta X_{k,\rm out}^{(i)},
 \qquad
 \gamma^{(i)}=(1-\eta)\gamma^{(i-1)}+\eta\gamma_{\rm out}^{(i)}.
 \label{eq:output-damping}
\end{equation}
Both first updates are used directly. This is arithmetic damping of the
message means and variances \cite[Sec.~V-B1]{cai2026}.

The stopping rule checks
$\norm{\widehat X^{(i)}-\widehat X^{(i-1)}}_F/
\max(\norm{\widehat X^{(i)}}_F,\norm{\widehat X^{(i-1)}}_F,10^{-12})$,
$\max_k|p_k^{(i)}-p_k^{(i-1)}|$, and
$|v^{(i)}-v^{(i-1)}|/\max(v^{(i)},v^{(i-1)})$.
Each must remain below a tolerance for a fixed number of iterations.
It also checks that the site update is accepted and that its precision
$\lambda_k=c_k^{-1}$, natural parameter $J_k=E_k/c_k$, and activity
history have stabilized. Specifically,
\begin{align}
 d_{\lambda,k}^{(i)}&=
 \frac{|\lambda_k^{(i)}-\lambda_k^{(i-1)}|}
      {\max(\lambda_k^{(i)},\lambda_k^{(i-1)})},\notag\\
 d_{J,k}^{(i)}&=
 \frac{\norm{J_k^{(i)}-J_k^{(i-1)}}}
      {\max(1,\norm{J_k^{(i)}},\norm{J_k^{(i-1)}})},\notag\\
 \max_k\{d_{\lambda,k}^{(i)},d_{J,k}^{(i)},
         |p_k^{(i)}-w_k^{(i-1)}|\}&<\epsilon_{\rm stop},
 \label{eq:site-stopping}
\end{align}
with $w_k^{(i)}=0.8w_k^{(i-1)}+0.2p_k^{(i)}$. The activity history is
used only for this check; the reported probabilities are $p_k^{(i)}$.
The two consecutive-iteration counters reset separately when their
respective conditions fail. The primary channel output is the
activity-weighted $\widehat X_k=p_k\mu_k$; the conditional estimate
$\mu_k$ is also retained, along with the decision in \eqref{eq:targets}.
Numerical stopping is not a global convergence guarantee.

\begin{center}
\begin{minipage}{0.95\linewidth}
\hrule\vspace{5pt}
\textbf{Algorithm 1: Channel and activity estimation}\par
\textbf{Input:} $\mathbf{Y},\mathbf{A},\sigma^2,\rho$ and trained networks
$S_\theta,T_\psi,z_\phi$.\par
\begin{enumerate}[leftmargin=1.8em,itemsep=2pt,topsep=4pt]
 \item Normalize by the fixed scale and initialize the prior and sites.
 Compute the thin SVD of $\mathbf{A}$.
 \item Compute the linear Gaussian mean, variance, and extrinsic query
 using \eqref{eq:linear-mean}--\eqref{eq:linear-extrinsic}; apply
 \eqref{eq:input-damping}.
 \item If $v$ is supported, query the mean and second networks once for
 each terminal, pool the variance by \eqref{eq:pooled-variance}, and
 update the sites by \eqref{eq:site-update} when $0<\tau<v$.
 Compute $p_k$ using the current classifier query
 \eqref{eq:activity-probability}. If $v$ is unsupported, retain the
 sites and use \eqref{eq:gaussian-evidence}.
 \item Form conditional moments, joint means, and spatial covariances
 using \eqref{eq:conditional-moments}--\eqref{eq:joint-spatial-covariance}.
 Compute and damp the outgoing message using
 \eqref{eq:joint-average-variance}--\eqref{eq:nonlinear-extrinsic} and
 \eqref{eq:output-damping}.
 \item Update the stopping statistics and activity history. Repeat
 steps 2--5 until the stopping rule or iteration limit is reached.
 \item Return $\widehat X_k^{\rm phys}=s\widehat X_k$, the conditional
 channel estimate $s\mu_k$, activity probabilities, decisions, and the
 stopping status. Rescale covariance outputs by $s^2$.
\end{enumerate}
\vspace{3pt}\hrule
\end{minipage}
\end{center}

The current implementation uses $M=16$, $N=16$, and $D=4096$.
The training bank contains 50,000 candidate channels from 12,500 scenes,
with four candidates per scene; validation uses 5,000 channels from
1,250 separate scenes. Query variances are sampled uniformly in $\log v$
over $[10^{-3},100]$. Active and inactive classifier queries share each
sampled variance and use independent Gaussian noise realizations.
The two channel networks pack the $M$ real and imaginary antenna maps
into $2M$ feature channels and use continuous log-variance conditioning.
Their base channel width is 128. The first network is frozen before
training the second network and classifier.

The fixed scale is $s=0.2647966453$, and the measured normalized training
power is $\widehat P_H=1.0004404181$. We use $\tau_{\min}=10^{-4}$,
$\varepsilon=10^{-6}$, $\eta=0.1$, and
$\epsilon_{\rm stop}=10^{-4}$. At least 60 iterations are performed;
both stopping checks must hold for ten consecutive iterations, with a
maximum of 600 iterations. A run reaching this maximum returns its
final estimate with an iteration-limit status.

With $Q\leq KN$, the one-time thin SVD costs $O(KNQ^2)$. Including the
implemented spatial transformations and small-matrix decompositions, one
supported iteration costs
$O(QKMN^2+KM^2N^2+M^3+K(C_\theta+C_\psi+C_\phi))$,
where the three $C$ terms are the per-terminal network forward costs.
Each network is evaluated once per terminal and iteration; the channel
mean is reused by the classifier. The Gaussian-evidence ablation omits
$C_\phi$, and unsupported queries omit all network forwards.
Excluding network parameters and activation storage, the main memory cost
is $O(QKN+KMN^2+KM^2)$. The present weights support the fixed
$M=N=16$ tensor size.

\begin{thebibliography}{8}
\bibitem{li2023}
K.-X. Li, X. Gao, and X.-G. Xia,
``Channel estimation for LEO satellite massive MIMO OFDM communications,''
\emph{IEEE Transactions on Wireless Communications}, 2023,
doi: \href{https://doi.org/10.1109/TWC.2023.3252895}{10.1109/TWC.2023.3252895}.
Author version: \url{https://arxiv.org/abs/2207.11958}.

\bibitem{huang2026}
Y. Huang, Y. Guo, C. Guo, Y. Wang, and X. Guo,
``Prior-agnostic AFDM channel estimation for high-mobility LEO satellite
communications using diffusion models,''
\emph{IEEE Wireless Communications Letters}, vol.~15,
pp.~2949--2952, 2026,
doi: \href{https://doi.org/10.1109/LWC.2026.3689601}{10.1109/LWC.2026.3689601}.

\bibitem{bemani2022}
A. Bemani, N. Ksairi, and M. Kountouris,
``Affine frequency division multiplexing for next generation wireless
communications,'' arXiv:2204.12798v2, 2022.
\url{https://arxiv.org/abs/2204.12798v2}.

\bibitem{cai2026}
C. Cai, W. Jiang, X. Yuan, and Y.-J. A. Zhang,
``Joint activity detection and channel estimation for massive connectivity:
Where message passing meets score-based generative priors,''
\emph{IEEE Journal on Selected Areas in Communications}, vol.~44, 2026,
doi: \href{https://doi.org/10.1109/JSAC.2026.3663274}{10.1109/JSAC.2026.3663274}.
Author version: \url{https://arxiv.org/abs/2506.00581}.

\bibitem{minka2001}
T. P. Minka, ``Expectation propagation for approximate Bayesian inference,''
in \emph{Proceedings of the Seventeenth Conference on Uncertainty in
Artificial Intelligence}, pp.~362--369, 2001.
\url{https://tminka.github.io/papers/ep/minka-ep-uai.pdf}.

\bibitem{efron2011}
B. Efron, ``Tweedie's formula and selection bias,''
\emph{Journal of the American Statistical Association},
vol.~106, no.~496, pp.~1602--1614, 2011,
doi: \href{https://doi.org/10.1198/jasa.2011.tm11181}{10.1198/jasa.2011.tm11181}.

\bibitem{vincent2011}
P. Vincent, ``A connection between score matching and denoising autoencoders,''
\emph{Neural Computation}, vol.~23, no.~7, pp.~1661--1674, 2011,
doi: \href{https://doi.org/10.1162/NECO_a_00142}{10.1162/NECO\_a\_00142}.

\bibitem{menon2016}
A. Menon and C. S. Ong,
``Linking losses for density ratio and class-probability estimation,''
in \emph{Proceedings of the 33rd International Conference on Machine Learning},
vol.~48, pp.~304--313, 2016.
\url{https://proceedings.mlr.press/v48/menon16.html}.
\end{thebibliography}
\end{document}
